% ----------------------
\subsection{Section 22 (16 April 1830)}
% ----------------------
	\label{DC22}
	
	% -----
	\subsubsection{Note on the JSP source}
	% -----
	
		(In case it isn't clear from what's below, this document does appear in RB1, but that isn't the JSP version. I have the RB1 version \hyperref[EXSS-DC22-RB1]{here} and note it in the verse comparisons below.)
		
		From the JSP: ``\textit{Revelation, [Fayette Township, Seneca Co., NY, 16 Apr. 1830]. Featured version part of `The Mormon Creed,' in} Painesville (OH) Telegraph\textit{, 19 Apr. 1831, vol. 2, no. 44 (second series), p. [4]. The microfilm copy of the text transcribed herein was filmed by the Microfilm Corporation of Cleveland, OH, 1947, copy at CHL.}
		
		``Both the \textit{Painesville Telegraph} version and the one found in Revelation Book 1 appear to have been created about the same time, but differences between the two versions indicate that the former was based on an earlier copy; therefore, the \textit{Telegraph} version is featured here.''
	
	% -----
	\subsubsection{Historical Background}
	% -----
		\label{DC22-HB}
		
		% --
		\paragraph{JSP}
		% --
		
			``JS dictated this revelation shortly after the formal organization of the church on 6 April 1830. While the version featured here does not include a specific date, a manuscript copy in the handwriting of William E. McLellin dates the revelation to 16 April 1830. When John Whitmer copied the revelation into Revelation Book 1, he likely wrote the heading found there: `A Revelation given to Joseph the Seer Some were anxious to Join the Church without Rebaptism \& Joseph enquired of the Lord \& he received as follows.'
			
			``Although several passages in the Book of Mormon emphasized the necessity of baptism by proper authority, no revelation prior to 16 April 1830 explicitly addressed the question of rebaptism for those who had been baptized in other faiths. Oliver Cowdery's June 1829 `Articles of the Church of Christ' prescribed the method of baptism and the wording of the baptismal prayer, declaring that `whosoever repenteth \& humbleth himself before me \& desireth to be baptized in my name shall ye baptize them,' but it did not address the question of rebaptism. The revelatory document on church government known as `Articles and Covenants,' which superseded Cowdery's earlier document, clarified that baptism was necessary for entry into the church but did not explicitly address rebaptism either.
			
			``The version of the 16 April revelation featured [in the JSP] was published in the \textit{Painesville Telegraph} and reportedly obtained from Martin Harris. The newspaper appended the revelation to its publication of Articles and Covenants, as though it were part of that text. Sidney Gilbert's early revelation notebook also appended the 16 April revelation to the end of Articles and Covenants, again without a separate heading or title. The first version of Articles and Covenants published in a church newspaper, in June 1832, likewise combined the two documents. The combining of these two texts in so many early versions suggests that the 16 April revelation was seen as an extension of the instructions on baptism contained in Articles and Covenants. The 16 April revelation is presented separately herein because the official version, found in Revelation Book 1, records it as a discrete text.
			
			``Despite the clarity this revelation may have provided church members, the requirement of rebaptism became a point of contention for many outside the church. Opponents criticized followers of JS for teaching `that their book contained a new covenant, to come under which the disciple must be re-immersed,' and when Thomas Campbell, father of Disciples of Christ founder Alexander Campbell, set forth his objections to JS's teachings, he argued that `re-baptizing believers is making void the ordinance of Christ.'''
			
	% -----
	\subsubsection{Versions}
	% -----
		\label{DC22-V}
		
		See the \href{https://drive.google.com/file/d/1goAvNz8jS4R8slYfMG_db55Ty2z_vmgK/view?usp=sharing}{sections comparison} document.
	
		\begin{compactdesc}
			\item[JSP] \hyperref[EMS-1830-JS-16-APR-1830]{16 April 1830}
			\item[RB1] \hyperref[EXSS-DC22-RB1]{here}
			\item[EMS] 1:1 (June 1832), 1
			\item[BC] Chapter 23, 47
			\item[DC 1835] Section 47, 178
			\item[TS] 4:1 (15 November 1842), 12
			\item[MS] 4:8 (December 1843), 116
			\item[DC 1844] Section 47, 267
		\end{compactdesc}

	% -----
	\subsubsection{Section 22:1} % *DC22-1 
	% -----
		\label{DC22-1}

		BEHOLD, I say unto you that all old covenants have I caused to be done away in this thing; and this is a new and an everlasting covenant, even that which was from the beginning.%
			\footnote{%
				% \label{}%
				So the ``new'' covenant is really the old \textit{old} covenant, the original one. Note, however, that this clarification is not in the JSP version, even though it is in RB1 (and it's actually even clearer in RB1).
				}

		\begin{framed}
		\scriptsize
			\begin{compactdesc}
				\item[JSP] A commandment unto the church of Christ which was established in these last days A. D. 1830, on the 4th month and the 6th day of the month, which is called April:--- Behold I say unto you that all old covenants have I caused to be done away in this thing, and this is a new and everlasting covenant;
				\item[RB1] A commandment unto the Church of Christ which was established in these \sout{day} last days one thousand eight hundred \& thirty on the forth month \& on the sixth day of the month which is called April Behold I say unto you that all old covenants have I \sout{called} caused to be done away in this thing \& this is a New \& an everlasting covenant even \sout{wherefore} the same which was from the begining
				% \item[BC] 
				% \item[]
			\end{compactdesc}
		\end{framed}

	% -----
	\subsubsection{Section 22:2} % *DC22-2 
	% -----
		\label{DC22-2}

		Wherefore, although a man should be baptized an hundred times it availeth him nothing, for you cannot enter in at the strait gate by the law of Moses, neither by your dead works.%
			\footnote{%
				% \label{}%
				By ``dead works'' I take the Lord to be referring to ordinances done without the full power of the priesthood.
				}

		\begin{framed}
		\scriptsize
			\begin{compactdesc}
				\item[JSP] wherefore, although a man should be baptized a hundred times it availeth him nothing, for ye cannot enter in at the straight gait by the law of Moses, neither by your dead works,
				\item[RB1] wherefore although a man shouldest be baptized an hundred times it availeth him nothing for ye cannot enter into the strait gate by the law of Moses neither by your dead works
				% \item[BC] 
				% \item[]
			\end{compactdesc}
		\end{framed}

	% -----
	\subsubsection{Section 22:3} % *DC22-3 
	% -----
		\label{DC22-3}

		For it is because of your dead works that I have caused this last covenant and this church to be built up unto me, even as in days of old.

		\begin{framed}
		\scriptsize
			\begin{compactdesc}
				\item[JSP] for it is because of your dead works that I have caused this last covenant and this church to be built up unto me;
				\item[RB1] for it is because of your dead works that I have caused this last covenant \& this church to be built up unto me even as in days of old
				% \item[BC] 
				% \item[]
			\end{compactdesc}
		\end{framed}

	% -----
	\subsubsection{Section 22:4} % *DC22-4 
	% -----
		\label{DC22-4}

		Wherefore, enter ye in at the gate, as I have commanded, and seek not to counsel your God.  Amen.

		\begin{framed}
		\scriptsize
			\begin{compactdesc}
				\item[JSP] wherefore enter ye in at the gate as I have commanded, and seek not to counsel your God.
				\item[RB1] wherefore enter ye in at the \sout{at} gate as I have commanded \& seek not to Council your God
				% \item[BC] 
				% \item[]
			\end{compactdesc}
		\end{framed}